%% Grandeurs de flux et grandeurs d'effort dans les quatre domaines (Tle, p. 245).
%% Un lien de puissance = un segment épais ; au-dessus le flux (vert), en dessous
%% l'effort (rouge) ; leur produit donne la puissance en watts.
\documentclass[tikz,border=4pt]{standalone}
\usepackage{liaisons}
\begin{document}
\begin{tikzpicture}[x=1cm,y=1cm,
  dom/.style={font=\scriptsize, anchor=east, align=right, text width=2.2cm, inner sep=2pt,
              execute at begin node={\hyphenpenalty=10000\exhyphenpenalty=10000}},
  lien/.style={line width=2pt, draw=black!65, line cap=round},
  flux/.style={font=\small, text=solideE, anchor=south, inner sep=3pt},
  eff/.style={font=\small, text=solideA, anchor=north, inner sep=3pt},
  puis/.style={draw=solideD, line width=1pt, fill=solideD!10, rounded corners=2pt,
               minimum width=2.05cm, minimum height=0.6cm, inner sep=2pt, font=\small}]

%% ---- macro : \ligne{y}{domaine}{flux}{effort}{puissance} -------------
\newcommand\ligne[5]{%
  \node[dom] at (2.3,#1) {#2};
  \draw[lien] (2.5,#1) -- (4.6,#1);
  \node[flux] at (3.55,#1) {#3};
  \node[eff]  at (3.55,#1) {#4};
  \node[puis] at (5.95,#1) {#5};}

\ligne{3.6}{Électrique (continu)}{$I$ (A)}{$U$ (V)}{$P = U \times I$}
\ligne{2.4}{Mécanique : translation}{$V$ (m$\cdot$s$^{-1}$)}{$F$ (N)}{$P = F \times V$}
\ligne{1.2}{Mécanique : rotation}{$\omega$ (rad$\cdot$s$^{-1}$)}{$C$ (N$\cdot$m)}{$P = C \times \omega$}
\ligne{0.0}{Hydraulique}{$Q$ (m$^3\cdot$s$^{-1}$)}{$p$ (Pa)}{$P = Q \times p$}

%% ---- rappels de couleur ---------------------------------------------
\node[font=\scriptsize, text=solideE, anchor=south] at (3.55,4.15) {grandeur de \textbf{flux}};
\node[font=\scriptsize, text=solideA, anchor=north] at (3.55,-0.5) {grandeur d'\textbf{effort}};

%% ---- conclusion ------------------------------------------------------
\node[font=\scriptsize, anchor=north] at (3.9,-0.95)
     {dans tous les domaines : \textbf{puissance = flux $\times$ effort}, en watts (W)};
\end{tikzpicture}
\end{document}
